% MAS1004 Assignment 2: long report.
%
% You fill this in with your agent, in your repository. AGENTS.md tells the
% agent how. Red text in the PDF is what still has to be replaced.
%
% - Pictures go in figures/. Copy them from results/, which is not in git.
% - Outputs go in outputs/: copy experiments.csv, LOG.md and the outputs of
%   check_submission.py and curate.py count here. \outputfile{outputs/...}
%   prints a file exactly as it is. Every number in this report has to come
%   from such an output.
%
% To make the PDF: compress this folder into a zip, and on Overleaf choose
% New Project, Upload Project. Set Menu, Compiler to XeLaTeX, and Recompile.

\documentclass[11pt]{article}
\usepackage{mas1004}

\reporttitle{Assignment 2: Long report}
\studentname{Your name}
\studentid{Your student ID}

\begin{document}
\makeheader

Leaderboard nickname: Lumi. Last releases: Track A \texttt{track-a-2}
(``every city drawn equally often'') and Track B \texttt{track-b-2}
(``MobileNetV3-Large at 512''). Every number below comes from
\texttt{outputs/experiments.csv}, \texttt{outputs/check\_submission.txt},
\texttt{outputs/leaderboard.txt} or the \texttt{count} outputs in
\texttt{outputs/}.

% ---------------------------------------------------------------------------
\section*{Track A}

All runs use seed 0 unless a seed is named. The starter recipe was trained
with seeds 0 and 1, which scored 0.4891 and 0.495: the seed alone moves the
score by about 0.006, so a change smaller than that may be chance.

\begin{center}\footnotesize
\begin{tabular}{lp{0.6\linewidth}r}
\toprule
tag & what changed from the run before & validation \\
\midrule
\texttt{baseline} (seed 0) & the starter recipe, nothing changed & 0.4891 \\
\texttt{baseline} (seed 1) & the same recipe, only the seed is different & 0.495 \\
\texttt{balanced} & \texttt{sample\_weights} gives every photo the weight one over the number of photos of its city, so every city is drawn about 20\% of the time & 0.5093 \\
\texttt{no\_flip} & \texttt{balanced} without \texttt{T.RandomHorizontalFlip()} & 0.4995 \\
\texttt{no\_panorama} & \texttt{balanced}, with the 670 training photos at least 2.0 times wider than tall left out of \texttt{selection.csv} & 0.5025 \\
\texttt{ema} & \texttt{balanced}, plus a running average of the weights (decay 0.999) that is handed in at the end & 0.509 \\
\bottomrule
\end{tabular}
\end{center}

\textbf{What I tried and why.} The reasons are in my own words in the lab
notebook at the end of this report.

\begin{itemize}
  \item \texttt{balanced}: ``because i remember in class the professor talked
  about the difference on the data each had, and because we cannot change the
  budget i thought about just having them equally.'' It scored 0.5093, the
  best Track A result, about 0.020 above baseline seed 0 and about 0.014
  above seed 1. By city it raised busan (0.3734 and 0.4 to 0.4532), osaka
  (0.3445 and 0.3769 to 0.4544) and taipei, and lowered tokyo (0.6254 and
  0.6035 to 0.5022). My reading: ``that the more data you have the better, so
  the better option would be to have more busan and osaka pictures instead of
  reducing and having them all equally but for now in track A, not
  possible.''
  \item \texttt{no\_flip}: ``i wanted to remove the flip because of text
  orientation and layour.'' It scored 0.4995, about 0.010 below
  \texttt{balanced}. My reading: ``probly not reading or not having legible
  hangul or kanji to identify places. because of the 224 by 224
  resolution.''
  \item \texttt{no\_panorama}: ``i say various of them on the suspicion
  sheets, thought it could be a way of cleaning data.'' It scored 0.5025,
  about 0.007 below \texttt{balanced}, close to the gap between the two
  baseline seeds. My reading: ``maybe it didn't work because the images were
  only like around 600? so the change wasnt very big.'' The 670 photos are
  1.5\% of the 44,257 training photos.
  \item \texttt{ema}: scored 0.509, the same as \texttt{balanced} (0.5093)
  to three decimals. The per-city scores moved in both directions.
\end{itemize}

\textbf{Final recipe and selection.} The final Track A release is
\texttt{track-a-2}: the \texttt{balanced} recipe, with no \texttt{selection.csv}
(all training photos). Its validation score on my run is 0.5093. On the
leaderboard page read on 2026-11-10 at 17:35 KST, \texttt{track-a-2} scored
0.414 (lowest and highest of the three seeds 0.398 and 0.436), with a
validation score of 0.504 on the grader's three runs. The instructor's
starter recipe scored 0.401 there. An earlier reading, at 15:06 KST, still
showed my starter recipe (score 0.399) because \texttt{track-a-2} had not been
scored yet. The rows are in \texttt{outputs/leaderboard.txt} below.

\outputfile{outputs/leaderboard.txt}

% ---------------------------------------------------------------------------
\section*{Decisions you made from photos}

\picturefile{figures/mistakes_baseline_seed0_busan_00.png}

\picturefile{figures/mistakes_baseline_seed0_tokyo_00.png}

\textbf{What I looked at and saw in common.} First the sheets of the
baseline's most confident mistakes (above, busan and tokyo), then the
sheets of the photos the \texttt{balanced} model doubts most (below). My notes
on the baseline's mistakes, in my own words:

\begin{quote}\small
I see a lot of cars without plates, and cities between same countries get
confused with each other when there is any writing or alphabets because its
the same, for example seoul or busan, so intra country regional confusion
cause it shares infrastructure or cultural elements. Buddhist temples are
similar between countries with lower resolution and a smaller image it could
probably not identify them correctly. big commercial places or markets can be
confusing too. generic asphalt is not as specific, food and indoor or close up
shots, panoramic or 360vies confuse the model too.
\end{quote}

\picturefile{figures/suspects_balanced_seed0_busan_00.png}

\textbf{The rule I wrote before acting.} ``lets drop photos with no
information, meaning photos that cannot be associated with the city. stuff
like indoors, food, extreme close ups, empty skies or stuff where nothing is
actually visible. so no information pictures.''

\textbf{What I changed.} I named 120 photos on the five suspects sheets
(busan 28, osaka 19, seoul 22, taipei 22, tokyo 29). They were taken out of
the selection with the reason ``no information: cannot be associated with the
city (indoors, food, extreme close ups, empty skies, nothing visible)'':

\outputfile{outputs/count_120_not_trained.txt}

That selection was prepared but never trained. I then chose a different
cleaning rule instead of it (``instead of cleaning only 120 pics''), and I
kept the 120 photos in for the \texttt{no\_panorama} run so that only one
thing changed. So the measured effect of the first rule is not known.

\textbf{The second rule, and the scores before and after.} I looked at two
sheets of wide photos (the 231 photos at least 2.5 times wider than tall, and
439 photos between 2.0 and 2.5) and decided: ``all the pics on the sheets look
like big panoramic photos. let the other pics stay in as i want to only change
1.'' The cut-off was width at least 2.0 times height, which left out 670
photos (the script is \texttt{scripts/make\_selection.py}):

\outputfile{outputs/count_no_panorama.txt}

Before: \texttt{balanced} 0.5093. After: \texttt{no\_panorama} 0.5025.

\picturefile{figures/wide_ratio_2p5_up_00.png}

\picturefile{figures/wide_ratio_2p0_to_2p5_00.png}

\textbf{The requests I gave my agent for these decisions, word for word.}

\begin{quote}\small
ok, lets drop photos with no information, meaning photos that cannot be
associated with the city. stuff like indoors, food, extreme close ups, empty
skies or stuff where nothing is actually visible. so no information pictures.
is this clear enough as a rule?
\end{quote}

\begin{quote}\small
for busan: 2, 4, 5, 6, 7, 9, 10, 13, 14, 15, 17, 18, 19, 23, 25, 27, 26, 28,
29, 30, 31, 32, 34, 35, 36, 38, 40, 43. for osaka: 1,3 , 5,7,9,12,14,17,18,19,26,
32, 35, 36, 37, 38, 42, 45, 47. for seoul: 1,2,4,5,9,11,13,14,17,19,26,27,28,29,
30,32,33,38,40,41,42,48. for taipei: 2,3,9,12,13,15,18,19,21,22,23,26,27,33,34,
37,40,41,43,44,45,46. for tokyo: 1,3,4,5,6,7,8,12,18,23,24,26,27,29,30,31,32,33,
34,35,36,40,41,42,43,44,45,46,48.
\end{quote}

\begin{quote}\small
now lets try a new change, cleaning rule. lets delete all the 260 or panoramic
like pictures. can we do that? instead of cleaning only 120 pics.
\end{quote}

\begin{quote}\small
all the pics on the sheets look like big panoramic photos. let the other pics
stay in as i want to only change 1.
\end{quote}

% ---------------------------------------------------------------------------
\section*{Track B}

Both runs trained on the starter training photos only (the \texttt{data}
line of \texttt{submission.json} says \texttt{data/starter/train}), 4 epochs.

\begin{center}\footnotesize
\begin{tabular}{lp{0.6\linewidth}r}
\toprule
tag & what it is & validation \\
\midrule
\texttt{b1} & \texttt{train\_b.py} as it came: ResNet18, 224 by 224 pixels & 0.4865 \\
\texttt{b2} & MobileNetV3-Large at 512 by 512 pixels, 4.20 million parameters, 1.11 billion multiply-adds & 0.5566 \\
\bottomrule
\end{tabular}
\end{center}

\textbf{What I tried.} In \texttt{b2} I changed two things at once, the model
and the image size, to fit the time I had on Colab, so the gain of about 0.070
cannot be split between them. ResNet18 cannot go above about 224 pixels within
the 2.0 billion multiply-adds limit, so a bigger image needed a smaller model.
Every city scored higher than in \texttt{b1}: seoul 0.639, busan 0.401, tokyo
0.644, osaka 0.440, taipei 0.659.
\textbf{My reason, in my own words:} ``when doing track A i found that letters
were probably not taken into account for now because of how small it is. the
224 x 224. because we cannot change for track A i wondered if we could for
Track B and because it was possible I did it. so there could be more info to
use and to find unique features from the similar cities.'' \textbf{What I
think the result says:} ``The model did well, but still is better on the main
cities, seoul and tokyo.''

\textbf{Extra photos.} I did not collect any: both Track B runs used only the
starter training photos.

\textbf{Validation next to leaderboard.} \texttt{b2} scored 0.5566 on the
validation photos. The leaderboard page (above, track B row) shows 0.432 for
\texttt{b2}. The leaderboard photos are frames from YouTube walking and driving
videos from channels that are not in the starter data, while the validation
photos are from the starter data.
\fillin{Why you think the two scores differ, in your own words.}

\subsection*{Your submission, as the grader sees it}
\outputfile{outputs/check_submission.txt}

% ---------------------------------------------------------------------------
\section*{What made the difference}

Track A: drawing every city equally often helped most (\texttt{balanced},
0.5093, against 0.4891 and 0.495 for the baseline). The other three changes
scored below it by about 0.010 (\texttt{no\_flip}), 0.007
(\texttt{no\_panorama}) and 0.0003 (\texttt{ema}); \texttt{ema} is far inside
the 0.006 gap between the two baseline seeds and \texttt{no\_panorama} is just
above it, so I cannot say either did anything.
Track B: the model and size change helped most, from 0.4865 to 0.5566. The
leaderboard scores of the starter recipe and \texttt{b2} are in
\texttt{outputs/leaderboard.txt}.

\fillin{What you think made the difference, in your own words.}

\fillin{What you would do next with one more week.}

% ---------------------------------------------------------------------------
\section*{All runs}
\outputfile{outputs/experiments.csv}

\section*{Lab notebook}
\outputfile{outputs/LOG.md}

\end{document}
